% ==================================================================
\section{Introduction}
\label{sec:intro}
% ==================================================================

Scientific publishing filters for success. Positive results
dominate the published literature across disciplines
\citep{fanelli2012negative}, and the pattern persists not because
scientists ask better questions but because negative findings are
systematically filtered out. Franco et al. found that roughly
two-thirds of null results from NSF-funded experiments were never
even written up \citep{franco2014publication}. The body of
scientific effort lost through this process is what Rosenthal
called the ``file drawer'' \citep{rosenthal1979file}.

The file drawer was not a failure of the system but an inevitable
triage under limited human retrieval and digestion bandwidth. When
the cost of finding a relevant failure exceeded the cost of simply
trying again, communities naturally prioritized sharing what worked
over documenting what did not.

This article argues that the arrival of large language models has
fundamentally altered this calculus, making the systematic
publishing of failure both feasible and urgent. We review the scale
and growing costs of the file drawer
(Section~\ref{sec:background}), analyze three ways in which LLMs
have changed the value of failure data
(Section~\ref{sec:opportunity}), propose experimental protocols for
validating these claims (Section~\ref{sec:future}), and discuss the
conditions for responsible use (Section~\ref{sec:discussion}).